\section{Related Work}

Forecasting plastic waste generation and other environmental indicators is often a small-data regression problem, where observations can be sparse across time, uneven across countries, and heterogeneous in measurement quality.

Environmental and waste-related forecasting studies increasingly evaluate classical time-series methods alongside supervised regression models that incorporate covariates, engineered temporal features, and regularization. Large comparative evaluations in environmental forecasting suggest that regression approaches (e.g., ridge regression and robust linear models) can be strong baselines when feature design is appropriate and overfitting is controlled \cite{effrosynidis2023time}. In waste-management contexts specifically, supervised learning models (including SVM/SVR variants) have been repeatedly explored for capturing nonlinear relationships between socioeconomic drivers and waste generation \cite{abbasi2013forecasting}.

% Forecasting plastic waste and related environmental indicators is often a \emph{small-data} regression problem: observations can be sparse across time, uneven across countries, and heterogeneous in measurement quality. These constraints tend to favor models that (i) work well with limited samples, (ii) regularize effectively to prevent overfitting, and (iii) support systematic hyperparameter selection to ensure fair model comparison. In this section, we review regression-based forecasting approaches relevant to sparse country-level time series, with emphasis on Support Vector Regression (SVR), Gaussian Process Regression (GPR), and hyperparameter optimization practices that align with the experimental design of this study.
%
% % \subsection{Regression Modeling for Sparse Environmental and Waste Forecasting}
%
% Environmental and waste-related forecasting studies increasingly evaluate classical time-series methods alongside supervised regression models that incorporate covariates, engineered temporal features, and regularization. Large comparative evaluations in environmental forecasting suggest that regression approaches (e.g., ridge regression and robust linear models) can be strong baselines when feature design is appropriate and overfitting is controlled \cite{Effrosynidis2023EnvForecastingComparison}. In waste-management contexts specifically, supervised learning models (including SVM/SVR variants) have been repeatedly explored for capturing nonlinear relationships between socioeconomic drivers and waste generation \cite{Abbasi2013MSW_SVR, Jassim2022SVR_MSW, Priyanka2025MSW_ML}.
%
% A key implication for plastic-waste forecasting is that model performance is often dominated not only by the learning algorithm but also by (1) preprocessing choices (missingness handling, scaling), (2) feature construction (lags, trends, per-capita transforms), and (3) careful selection of model capacity via regularization and hyperparameters. As a result, modern forecasting pipelines frequently treat “model selection” as joint optimization over the full pipeline rather than the estimator alone.
%
% % \subsection{Support Vector Regression for Small-Sample Nonlinear Forecasting}
%
% SVR is a kernel-based regression method grounded in statistical learning theory \cite{Vapnik1995NatureSLT}. It fits a function by minimizing a regularized objective with an $\epsilon$-insensitive loss, balancing empirical error against model complexity. This bias--variance control is particularly attractive for sparse, noisy datasets because capacity can be adjusted via the regularization parameter $C$, the margin parameter $\epsilon$, and kernel-specific parameters (e.g., $\gamma$ for the RBF kernel) \cite{Smola2004SVRTutorial}.
%
% In environmental and waste-related applications, SVR is commonly used to model nonlinear effects from demographic and economic covariates and to stabilize predictions under limited data. Prior work on municipal solid waste forecasting reports competitive performance for SVR-like approaches, often outperforming simpler linear baselines when relationships are nonlinear and interactions matter \cite{Abbasi2013MSW_SVR, Jassim2022SVR_MSW}. However, SVR performance is notably sensitive to scaling, kernel choice, and hyperparameter settings, which makes systematic tuning essential for fair evaluation across countries and metrics.
%
% % \subsection{Gaussian Process Regression and Uncertainty-Aware Forecasts}
%
% GPR provides a Bayesian, nonparametric approach to regression by placing a Gaussian process prior over functions \cite{Rasmussen2006GPML}. Predictions are probabilistic, producing both a mean estimate and an uncertainty measure that reflects data scarcity, noise assumptions, and kernel structure. This property is valuable in sparse forecasting settings because it helps distinguish confident extrapolations from regions where data support is weak.
%
% Kernel design is central in GPR: covariance functions encode assumptions about smoothness, periodicity, and trends. For time-series and environmental data, GPR has been used to capture structured temporal patterns and to incorporate domain knowledge through composite kernels \cite{Roberts2013GPTimeSeries, WilliamsRasmussen1996GPR}. In addition, uncertainty-aware modeling is often highlighted as a practical advantage for decision-making systems that must communicate prediction reliability, especially when country-level coverage is uneven or missing.
%
% Despite these advantages, vanilla GPR has computational limits ($\mathcal{O}(n^3)$ training in the number of observations) and is sensitive to kernel and noise hyperparameters; consequently, robust experimentation typically requires well-defined tuning procedures and consistent evaluation protocols.
%
% % \subsection{Hierarchical and Grouped Forecasting Considerations}
%
% When forecasts are needed at multiple aggregation levels (e.g., country, region, global), hierarchical time-series methods aim to reconcile predictions so that totals and subseries remain coherent. Forecast reconciliation methods formalize this constraint, producing adjusted forecasts that respect aggregation structure while minimizing error \cite{Wickramasuriya2019ForecastReconciliation}. While this work primarily emphasizes country-level regression forecasting, hierarchical considerations remain relevant if subsequent extensions require consistent roll-ups across geographic groupings or product categories.
%
% % \subsection{Hyperparameter Optimization and Fair Model Comparison}
%
% Because SVR and GPR are highly hyperparameter-dependent, the credibility of comparative results hinges on tuning methodology and validation design. Common approaches include grid search and random search; random search is often preferred in high-dimensional spaces because it explores more diverse configurations under the same budget \cite{BergstraBengio2012RandomSearch}. Beyond these baselines, Bayesian optimization methods model validation performance as a function of hyperparameters and select new trials using an acquisition rule, often yielding strong results with fewer evaluations \cite{Snoek2012PracticalBO, Hutter2011SMAC}. Bandit-based methods such as Hyperband further reduce cost by early-stopping poorly performing configurations \cite{Li2017Hyperband}.
%
% Equally important is the validation protocol. Cross-validation remains a standard tool for estimating generalization and supporting model selection under limited data \cite{Kohavi1995CrossValidation}. In this study, the optimization loop is designed to tune hyperparameters by minimizing MAE under $k$-fold cross-validation, consistent with established practice for small-sample regression and aligned with the later experimental plan. This alignment is also essential for comparing SVR, GPR, and other regression baselines under a shared objective and budget (e.g., multi-start strategies), helping ensure that observed differences reflect modeling assumptions rather than uneven tuning effort.
%
% % \subsection{Summary of Gaps and Motivation}
%
% Across the literature, two recurring gaps appear in sparse environmental/waste forecasting: (1) comparisons that do not control for tuning effort across models, and (2) pipelines that do not explicitly connect preprocessing, feature design, and hyperparameter selection to a unified evaluation protocol. These issues can lead to misleading conclusions, especially when data availability varies substantially across countries.
%
% Motivated by these gaps, this work emphasizes regression methods suited to small datasets (notably SVR and GPR) and adopts a consistent hyperparameter optimization strategy using cross-validated MAE to support fair, reproducible comparison across forecasting targets and geographic subsets.
